\documentclass[11pt]{letter}
\usepackage[margin=1in]{geometry}
\usepackage{hyperref}

\signature{Liang Wang, Ph.D.\\
School of Artificial Intelligence and Automation\\
Huazhong University of Science and Technology}

\address{Liang Wang\\
School of Artificial Intelligence and Automation\\
Huazhong University of Science and Technology\\
Wuhan, Hubei 430070, P.R. China\\
Email: wangliang.f@gmail.com}

\begin{document}

\begin{letter}{Editor-in-Chief\\
BMC Bioinformatics\\
(or PLOS Computational Biology)}

\opening{Dear Editor,}

I am writing to submit our manuscript titled \textbf{``OmniGene-4: Mixture-of-Experts Routing Reveals a 96\%/4\% CPT/SFT Decomposition in Biological Foundation Models''} for consideration as a Research Article in \emph{BMC Bioinformatics} (or \emph{PLOS Computational Biology}).

\section*{Significance and Novelty}

Mixture-of-Experts (MoE) architectures are increasingly adopted in large language models, yet their internal routing mechanisms remain poorly understood, particularly in biological applications. Our work makes three novel contributions:

\begin{enumerate}
\item \textbf{First systematic MoE routing analysis in biological foundation modeling.} We introduce a forward-hook-based framework to collect token-level expert activations across eight biological task families (DNA, protein, structure, cell, molecule, literature), revealing task-specific expert specialization that emerges during training.

\item \textbf{Quantitative CPT/SFT decomposition with statistical rigor.} Through prompt-level bootstrap resampling (1,000 iterations), we demonstrate that continued pretraining (CPT) accounts for 96.2\% [95\% CI: 95.9\%, 96.6\%] of cross-task routing differentiation, while supervised fine-tuning (SFT) contributes only 3.7\% [3.4\%, 4.1\%]. This is the first such measurement reported for a biological MoE and challenges the conventional wisdom that task-specific fine-tuning is the primary driver of model specialization.

\item \textbf{Robustness validation through format control and external benchmarking.} We address two critical confounds: (a) format-matched routing control confirms 90.2\% of routing signal is content-driven rather than prompt-format artifact; (b) head-to-head comparison on identical 2,000-pair remote homology evaluation shows OmniGene-4 (59.5\%) outperforms ESM-2 650M (50.5\%) by 9 percentage points, demonstrating that instruction-tuned general LLMs with biological CPT can match or exceed specialist protein-embedding models.
\end{enumerate}

\section*{Methodological Rigor}

Our work adheres to high standards of reproducibility and statistical transparency:
\begin{itemize}
\item All training, evaluation, and analysis scripts are publicly available in the project repository (\url{https://github.com/maris205/omnigene4}).
\item Bootstrap confidence intervals quantify uncertainty in all routing measurements.
\item Limitations are explicitly disclosed, including single-model-family constraint, modest remote-homology ceiling, and benchmark ownership.
\end{itemize}

\section*{Target Audience and Impact}

This work is directly relevant to the bioinformatics community for three reasons:
\begin{enumerate}
\item \textbf{Practitioners} building biological foundation models can use our routing-analysis framework to audit expert specialization and diagnose catastrophic forgetting.
\item \textbf{Researchers} studying model interpretability gain a discrete, auditable observable (router activations) that complements continuous probes like sparse autoencoders.
\item \textbf{Computational biologists} seeking to understand when CPT vs.\ SFT is more effective can use our decomposition as a guide for resource allocation.
\end{enumerate}

\section*{Why This Journal}

We believe \emph{BMC Bioinformatics} (or \emph{PLOS Computational Biology}) is the ideal venue for this work because:
\begin{itemize}
\item The journal values methodological innovation and reproducibility, which are central to our contribution.
\item The readership includes both machine learning practitioners and computational biologists, matching our interdisciplinary approach.
\item The open-access model aligns with our commitment to making biological foundation models accessible to the broader research community.
\end{itemize}

\section*{Suggested Reviewers}

We respectfully suggest the following experts as potential reviewers:

\begin{enumerate}
\item \textbf{Dr.\ Zeming Lin} (Meta AI / MIT)\\
Email: zeminl@mit.edu\\
Expertise: ESM-2, protein language models\\
Reason: Co-author of ESM-2, directly relevant to our external benchmarking.

\item \textbf{Dr.\ Sergey Ovchinnikov} (MIT)\\
Email: so@mit.edu\\
Expertise: Protein structure prediction, deep learning\\
Reason: Expert in protein foundation models and interpretability.

\item \textbf{Dr.\ Jian Peng} (University of Illinois Urbana-Champaign)\\
Email: jianpeng@illinois.edu\\
Expertise: Computational biology, deep learning\\
Reason: Broad expertise in biological sequence modeling.

\item \textbf{Dr.\ Jinbo Xu} (Toyota Technological Institute at Chicago)\\
Email: jinboxu@gmail.com\\
Expertise: Protein structure prediction, machine learning\\
Reason: Expert in protein homology detection and deep learning methods.
\end{enumerate}

\section*{Conflicts of Interest}

The author declares no conflicts of interest. The BioPAWS benchmark used in this study was developed by the author, which is disclosed in the Limitations section of the manuscript.

\section*{Data and Code Availability}

All code, training scripts, and analysis notebooks are available at \url{https://github.com/maris205/omnigene4}. The routing-count matrices and token-level activation database are in the \texttt{outputs/moe\_analysis/} directory. LoRA adapter weights ($\approx 1.9$ GB) can be released upon reasonable request.

\vspace{1em}

Thank you for considering our manuscript. We look forward to your response.

\closing{Sincerely,}

\end{letter}

\end{document}
